%% Q1
A \textbf{semantic barrier} arises from problems of \emph{meaning}: words, symbols, jargon or technical terms that the receiver does not understand, or understands differently. Using technical or difficult language is the classic example.
\why{Use of technology relates to channel or technological barriers; lack of attention is a psychological barrier; lack of co-ordination is an organisational barrier.}

%% Q2
Communication is purposeful: it aims to inform, persuade, motivate and so \textbf{bring about a change} in the knowledge, attitude or behaviour of the receiver. Effective communication achieves the intended effect.
\why{Barriers, misunderstanding and confusion are the \emph{opposite} of effective communication.}

%% Q3
In \textbf{Aristotle's model} of communication (the earliest linear model) the elements occur in this order:
\begin{center}
\begin{tikzpicture}[node distance=6mm,every node/.style={draw=brand,rounded corners,fill=soft,font=\small,minimum height=6mm}]
\node (a) {Speaker (b)}; \node[right=of a] (b) {Speech (a)}; \node[right=of b] (c) {Receiver (c)}; \node[right=of c] (d) {Effect (d)};
\draw[-Stealth,thick,brand] (a)--(b); \draw[-Stealth,thick,brand] (b)--(c); \draw[-Stealth,thick,brand] (c)--(d);
\end{tikzpicture}
\end{center}
Speaker $\rightarrow$ Speech $\rightarrow$ Audience (receiver) $\rightarrow$ Effect, i.e.\ \textbf{b, a, c, d}. (Lasswell's later model follows the same logic: \emph{Who says What, in Which channel, to Whom, with What effect}.)

%% Q4
\textbf{Non-verbal} communication uses no words: gestures, facial expressions, posture, eye contact, head movements, touch, space and time. \textbf{Nodding the head} (to agree or disagree) is a gesture, which falls under kinesics.
\why{Writing on the blackboard is written verbal communication; speaking on a microphone and sending a voice message are oral verbal communication.}

%% Q5
\textbf{Encoding} is the sender's act of converting thoughts, ideas or feelings into a message using symbols (words, pictures, gestures).
\why{Option (B) describes \emph{decoding}, which is done by the receiver.}

%% Q6
\textbf{Python} is the most widely used language in present-day AI and machine learning, thanks to libraries such as TensorFlow, PyTorch, scikit-learn, NumPy and Keras. The key treats it as the common language for AI.
\why{LISP (1958) was the historical language of AI research, but today it is niche. PHP is a web scripting language, and Java, though usable, is not the common AI language.}

%% Q7
All three statements are correct:
\begin{itemize}
\item ICT increases \textbf{transparency} (online records, RTI portals, e-tendering).
\item e-Governance involves \textbf{automation of administrative processes} (e-office, online services).
\item Automation reduces paperwork and manual jobs, so the \textbf{workforce needed is reduced}.
\end{itemize}
Hence i, ii and iii.

%% Q8
SWAYAM has \textbf{nine National Coordinators}, each responsible for one type of content: AICTE (self-paced and international courses), \textbf{NPTEL/IIT Madras} (engineering), \textbf{UGC} (non-technical postgraduate education), \textbf{CEC} (undergraduate education), \textbf{NCERT} and \textbf{NIOS} (school education), IGNOU (out-of-school students), \textbf{IIM Bangalore} (management studies) and NITTTR (teacher training).
Among the given names, UGC, CEC, IIMB and NIOS are national coordinators. IIT Kanpur is not (NPTEL is coordinated through IIT Madras), so the answer is \textbf{2, 3, 4 and 5}.

%% Q9
\textbf{Statement I} is true: SWAYAM (Study Webs of Active-Learning for Young Aspiring Minds), launched in 2017 by MHRD, is India's own MOOC platform offering free courses from Class~9 to postgraduate level. \textbf{Statement II} is true: Coursera, Udemy, Udacity and iversity are well-known international MOOC providers.

%% Q10
A touch screen both \textbf{displays} output and \textbf{accepts} touch input, so it is an input as well as an output device. The reason, that it can send and receive data at the same time, correctly explains why it plays both roles. Both are true and (R) explains (A).

%% Q11
In Dunkin and Biddle's \textbf{Presage--Process--Product} model:
\begin{itemize}
\item \textbf{Presage} variables come \emph{before} teaching: teacher characteristics, planning, \emph{content analysis}, the \emph{learning outcome plan}, and learner characteristics such as \emph{motivation}.
\item \textbf{Process} variables operate \emph{during} teaching: classroom interaction and learning supervision.
\item \textbf{Product} variables are the \emph{results}: learning and its testing.
\end{itemize}
Presage components here are 1, 4 and 6.

%% Q12
The discussion method promotes understanding, reflective thinking and active participation, but because it follows the flow of the group's ideas it does \textbf{not} guarantee that content is organised in a structured, logical sequence. That is the strength of the lecture or programmed method.

%% Q13
The teacher observed that brief summaries help learners \emph{memorise and recall}. Sensible implications are that summaries motivate participation, aid understanding and conceptualisation, and support cognitive learning. Saying that summarisation focuses on memorisation \textbf{irrespective of understanding} contradicts good pedagogy and is \emph{not} a valid implication.

%% Q14
Online teaching allows rich \textbf{multimedia} (video, animation, simulation), requires content to be \textbf{well structured} into modules, and is built around clear \textbf{learning objectives}. All three are advantages, so the answer is \emph{all of the above}.

%% Q15
Factors that \emph{disturb} classroom teaching:
\begin{itemize}
\item 2. absence of teacher clarification, which leaves learners confused;
\item 4. classroom noise by learners;
\item 5. peer attention-seeking arguments.
\end{itemize}
Self-confidence (1) and logical arguments for clarification (3) help learning, and access to digital material (6) is a resource rather than a disturbing \emph{characteristic}. The answer is 2, 4, 5.

%% Q16
Rule: each term $=\dfrac{\text{previous}}{2}-k$, with $k$ decreasing by 2:
\[864\xrightarrow{\div2-12}420\xrightarrow{\div2-10}200\xrightarrow{\div2-8}\mathbf{92}\xrightarrow{\div2-6}40\xrightarrow{\div2-4}16\xrightarrow{\div2-2}6\]
The series should have \textbf{92}, so \textbf{96} is the wrong number (and $92/2-6=40$ confirms it).

%% Q17
Semi-annual: rate $=5\%$ per half-year, $n=3$ half-years.
\[CI=P[(1.05)^3-1]=0.157625P,\qquad SI=P\times10\%\times1.5=0.15P\]
\[CI-SI=0.007625P=183\ \Rightarrow\ P=\frac{183}{0.007625}=\textbf{Rs.\ 24{,}000}\]

%% Q18 fig
\small
\begin{tabular}{c c cc}\toprule
& & \multicolumn{2}{c}{\textbf{Males : Females}}\\\cmidrule(l){3-4}
\textbf{State} & \textbf{\% below poverty line} & \textbf{Below PL} & \textbf{Above PL}\\\midrule
A & 35 & 5 : 6 & 6 : 7\\ B & 25 & 3 : 5 & 4 : 5\\ C & 24 & 1 : 2 & 2 : 3\\ D & 19 & 3 : 2 & 4 : 3\\ E & 15 & 5 : 3 & 3 : 2\\\bottomrule
\end{tabular}

%% Q18 after
What will be the number of females above poverty line in the State D if it is known that the population of State D is 7 million?

%% Q18 q
The following table gives the percentage distribution of population of five States A, B, C, D and E on the basis of poverty line and also on the basis of gender:

%% Q18
State D: 19\% below poverty line, so $100-19=81\%$ are above.\\
Population above poverty line $=81\%\times7=5.67$ million.\\
Males : Females above PL $=4:3$, so females $=\tfrac37\times5.67=\textbf{2.43 million}$.

%% Q19
Cricket only $+$ Football only $+$ Both $=5+15+75=95\%$, so non-members $=5\%$.
\[5\%\text{ of total}=40\ \Rightarrow\ \text{total}=800\qquad \text{Both clubs}=75\%\times800=\mathbf{600}\]

%% Q20 fig
\small
\begin{tabular}{lccccc}\toprule
\textbf{Day} $\backslash$ \textbf{Machine} & \textbf{A} & \textbf{B} & \textbf{C} & \textbf{D} & \textbf{E}\\\midrule
Monday & 139 & 147 & 211 & 141 & 184\\ Tuesday & 141 & 189 & 164 & 189 & 151\\ Wednesday & 115 & 141 & 159 & 156 & 136\\ Thursday & 89 & 223 & 120 & 147 & 113\\ Friday & 187 & 93 & 257 & 160 & 124\\\bottomrule
\end{tabular}

%% Q20 after
The number of auto-parts manufactured by Machine E on Monday is what percent more than the number of parts manufactured by Machine C on Thursday?

%% Q20 q
The number of auto-parts manufactured by 5 identical machines during 5 days is given below:

%% Q20
Machine E on Monday $=184$; Machine C on Thursday $=120$.
\[\frac{184-120}{120}\times100=\frac{64}{120}\times100=\mathbf{53.33\%}\]

%% Q21
Since 2019 \textbf{NEET-UG} has been conducted by the \textbf{National Testing Agency (NTA)}, an autonomous body set up in 2017 under the Ministry of Education. It also conducts JEE (Main), CUET and UGC-NET.
\why{CBSE conducted NEET before 2019. The Medical Council of India was replaced by the National Medical Commission in 2020.}

%% Q22
\textbf{ICSSR} is the \textbf{Indian Council of Social Science Research}, established in 1969 under the Ministry of Education to promote research in the social sciences.
\tip{Other councils: ICHR (Historical Research, 1972), ICPR (Philosophical Research, 1977), ICAR (Agricultural Research), ICMR (Medical Research), CSIR (Scientific \& Industrial Research).}

%% Q23
In \textbf{NIRF 2024} (India Rankings, Ministry of Education), \textbf{IIT Madras} was ranked No.~1 in the \emph{Overall} category for the sixth consecutive year (2019--2024). IISc Bangalore was 2nd and IIT Bombay 3rd.
\tip{In the \emph{University} category, IISc Bangalore is ranked 1st.}

%% Q24
\textbf{NIRF} is the \textbf{National Institutional Ranking Framework}, launched by MHRD in 2015 (first rankings 2016). Its five parameters are Teaching, Learning \& Resources; Research \& Professional Practice; Graduation Outcomes; Outreach \& Inclusivity; and Perception.

%% Q25
The \textbf{Malaviya Mission Teacher Training Programme (MMTTP)} was launched by the Ministry of Education in 2023, renaming and revamping the earlier Human Resource Development Centres (formerly Academic Staff Colleges). Its \textbf{Malaviya Mission Teacher Training Centres} build the competencies of higher-education faculty in line with NEP~2020.

%% Q26
Basic (fundamental) research aims at \emph{general, objective} knowledge. Its findings are replicable, its claims generalisable, and its concepts and procedures publicly accessible for verification. Being \textbf{human-centric} (value-laden, focused on human problems) is not a defining characteristic of basic research; it describes applied or action-oriented research better.

%% Q27
\textbf{MS Word} is a word processor for creating documents; it is not a data-analysis tool.
\why{MS Excel (functions, pivot tables, Data Analysis ToolPak), SPSS (Statistical Package for the Social Sciences) and MATLAB are all used for quantitative analysis.}

%% Q28
The researcher wants to \emph{explore} the factors behind organising a support system for girl students as it exists or could exist on campus, typically by gathering views from students, teachers and administrators. Collecting such present-status information from a group is the \textbf{descriptive survey method}.
\why{Historical research studies past events; philosophical research analyses concepts and values; experimental research manipulates variables under control, which is not the aim here.}

%% Q29
In a normal distribution: $\pm1\sigma$ covers 68.26\%, $\pm2\sigma$ covers \textbf{95.44\%} ($\approx$ 95\%), and $\pm3\sigma$ covers 99.73\%.

%% Q30
In \textbf{non-probability sampling} not every unit has a known chance of selection. Examples are quota, purposive, convenience, snowball and judgement sampling. \textbf{Quota sampling} is non-probabilistic.
\why{Cluster, systematic and simple random sampling are all \emph{probability} sampling techniques.}

%% Q31
\[\frac{(W-2)^2}{V}=\frac{(3-2)^2}{1/2}=2,\qquad \frac{(Y-5)^2}{Z}=\frac{(9-5)^2}{1/6}=96,\qquad U-X=-2-(-4)=2\]
\[\sqrt{2+96+2}=\sqrt{100}=\mathbf{10}\]

%% Q32
$16\tfrac23\%=\tfrac16$. The CI factor for 3 years is $\left(\tfrac76\right)^3-1=\tfrac{343}{216}-1=\tfrac{127}{216}$.
\[P=1270\times\frac{216}{127}=2160\qquad SI=2160\times\frac16\times3=\textbf{Rs.\ 1{,}080}\]

%% Q33
Filling in \textbf{b, a, b, b, a} gives
\[\text{ab\underline{b}a}\ \ \text{\underline{a}b\underline{b}a}\ \ \text{ab\underline{b}a}\ \ \text{abb\underline{a}}\]
The block ``abba'' repeats four times, so the answer is ``babba'', option (B).

%% Q34
Let the second number be $6k$. Then the first is $3k$ and the third is $2k$.
\[3k+6k+2k=3\times44=132\Rightarrow 11k=132\Rightarrow k=12\]
The numbers are 36, 72 and 24, so the largest is \textbf{72}.

%% Q35
After 10\% die, 90\% remain; then 20\% of these leave, so 80\% of 90\% remain.
\[P\times0.9\times0.8=5040\Rightarrow 0.72P=5040\Rightarrow P=\mathbf{7000}\]

%% Q36
SDG~4 Quality education (c); SDG~6 Clean water \& sanitation (a); SDG~7 Affordable \& clean energy (d); SDG~13 Climate action (b). The match is \textbf{I-c, II-a, III-d, IV-b}.

%% Q37
The key treats both statements as true. Earthquakes, tsunamis, floods and avalanches are standard examples of \textbf{natural hazards} (Statement II). Statement I describes the disasters referred to as arising from natural hazardous processes, and the official key accepts it as true.
\tip{Hazard = a potentially damaging natural process. Disaster = a hazard that actually causes serious loss because the community is vulnerable.}

%% Q38
Global warming leads to changed rainfall patterns, sea-level rise (thermal expansion and melting ice), loss of biodiversity, and shrinking forests (droughts, forest fires, species shifts). All four are effects, so the answer is (C).

%% Q39
West coast (Arabian Sea) cyclones: \textbf{Biparjoy} (June 2023, landfall near Jakhau, Kutch, Gujarat) and \textbf{Tauktae} (May 2021, landfall in Saurashtra, Gujarat).
East coast (Bay of Bengal) cyclones: \textbf{Fani} (2019, Odisha) and \textbf{Amphan} (2020, West Bengal).
The answer is 1 and 3.

%% Q40
Soil is polluted by pesticides and fertilisers, radioactive waste, municipal solid waste, mining waste (heavy metals, acid mine drainage) and e-waste (lead, mercury, cadmium). All five are soil pollutants.

%% Q41 ctx
\textbf{Directions for Q.~41 to 45:} Read the following paragraph and answer the questions.\par\smallskip
There are several other fields in which it is desirable that a clear policy should be laid down on an all-India basis. One of the fundamental rights guaranteed by the Constitution is the right of every citizen of the Union to free and compulsory education up to the age of 14. For the proper functioning of democracy, the Centre must see that every individual is equipped with the necessary knowledge, skill and aptitudes to discharge his duties as a responsible and co-operative citizen. What James Madison said about his country, the United States, many years ago, holds good today in our country also. ``A popular Government without popular information or the means of acquiring it is but a prologue to a farce of tragedy or perhaps both. Knowledge will forever govern ignorance, and the people who mean to be their own governors must arm themselves with the power which knowledge gives.'' We should, however like to add that it is not only knowledge that is required, but also the right kind of social training and the inculcation of right ideals without which knowledge by itself may be sterile or worse. Training for democracy postulates a balanced education in which social virtues, intellectual development and practical skill all receive due consideration and the pattern of such an education must be envisaged on an all-India basis.

%% Q41
The passage states: ``One of the fundamental rights guaranteed by the Constitution is the right of every citizen \ldots\ to \textbf{free and compulsory education} up to the age of 14.'' (Article 21A was added by the 86th Amendment, 2002, and implemented through the RTE Act, 2009.)

%% Q42
``\ldots every individual is equipped with the necessary knowledge, skill and aptitudes \textbf{to discharge his duties as a responsible and co-operative citizen}.''

%% Q43
Madison said that a popular government without popular information ``is but a \textbf{prologue to a farce of tragedy}.'' Option (C) paraphrases this.

%% Q44
``Training for democracy postulates a \textbf{balanced education} in which \textbf{social virtues, intellectual development and practical skill} all receive due consideration.''

%% Q45
The author wants a policy on an all-India basis so that every citizen receives the same kind of \textbf{balanced education}, with uniform standards across the country. The pattern ``must be envisaged on an all-India basis.''

%% Q46
\begin{center}
\begin{tikzpicture}[scale=0.9,every node/.style={font=\small}]
\draw[->,gray] (-0.3,-0.3)--(-0.3,0.5) node[above]{N};
\fill (0,0) circle(2pt) node[below left]{Z};
\fill (0,2) circle(2pt) node[above left]{Y};
\fill (2.5,2) circle(2pt) node[above right]{X};
\fill (2.5,0) circle(2pt) node[below right]{P};
\draw[thick,brand] (0,0)--(0,2)--(2.5,2); \draw[thick,brand] (0,0)--(2.5,0);
\draw[-Stealth,thick,accent,dashed] (2.5,2)--(2.5,0.1);
\end{tikzpicture}
\end{center}
Y is north of Z. X is east of Y. P is east of Z. Taking the standard layout, X and P are at the same east--west position, with X directly above P. So P is to the \textbf{South} of X.

%% Q47
In Nyaya epistemology:
\begin{itemize}
\item \textbf{Pramana}: the means of valid knowledge (b)
\item \textbf{Pramata}: the knower, the subject who knows (a)
\item \textbf{Prameya}: the object of knowledge (d)
\item \textbf{Prama}: the resulting valid knowledge (c)
\end{itemize}
The match is \textbf{1-b, 2-a, 3-d, 4-c}.

%% Q48
New ratio $=6\times1.20 : 8\times1.15 : 7\times1.20=7.2:9.2:8.4=72:92:84=\mathbf{18:23:21}$.

%% Q49
\[\frac{1+4+9+16+25+36}{6}=\frac{91}{6}=\mathbf{15.17}\]
(Using $\sum n^2=\tfrac{n(n+1)(2n+1)}{6}=\tfrac{6\cdot7\cdot13}{6}=91$.)

%% Q50
The growth is compound (per day): $1000(1+r)^3=4000\Rightarrow(1+r)^3=4\Rightarrow1+r=\sqrt[3]4\approx1.587$.
Average daily increase $\approx\mathbf{58.7\%}$.
